\section{总结与展望}
\subsection{感想总结}
\begin{table}[htbp]
    \centering
    \begin{tabular}{|l|l|}
    \hline
    \multicolumn{1}{|c|}{{\color[HTML]{000000} \textbf{毕业要求}}}                             & \multicolumn{1}{c|}{{\color[HTML]{000000} \textbf{课程对应毕业要求指标点}}}                                                                                      \\ \hline
    {\color[HTML]{000000} \textbf{\begin{tabular}[c]{@{}l@{}}3.设计/\\ 开发解决方案\end{tabular}}} & \cellcolor[HTML]{FFFFFF}{\color[HTML]{000000} \begin{tabular}[c]{@{}l@{}}3-2 能够针对特定指标需求，设计并实现功能完整的硬件和软件\\ 系统，包括整体架构设计、各模块交互和数据通信等（系统设计）\end{tabular}} \\ \hline
    {\color[HTML]{000000} \textbf{4. 研究}}                                                  & \cellcolor[HTML]{FFFFFF}{\color[HTML]{000000} \begin{tabular}[c]{@{}l@{}}4-2 能够根据方案，运用实验工具、仪器，开展实验，对实验\\ 结果进行分析与解释，并通过信息综合得到合理有效的结论\end{tabular}}     \\ \hline
    {\color[HTML]{000000} \textbf{10.沟通}}                                                  & {\color[HTML]{000000} \begin{tabular}[c]{@{}l@{}}10-1能够就电子信息工程领域的复杂工程问题撰写报告，设计文稿，\\ 并具有较好的语言表达和沟通能力，能够清晰陈述观点和回答问题\end{tabular}}                       \\ \hline
    \end{tabular}
    \end{table}
在参与整个项目的过程中，我深刻体会到了团队合作、技术创新以及工程实践的重要性。作为一项综合性极强的嵌入式系统设计扩充，该项目不仅需要扎实的理论基础，还需要高效的实践和调试能力。在设计与制作过程中，我学到了如何将传感器、嵌入式控制、电机驱动等技术结合在一起，解决实际问题。

在构建项目的过程中，我们面临了不少挑战。首先是如何平衡系统的稳定和速度。为了实现这一目标，我们反复调试气垫系统的气压分布以及行进电机的角度和功率，通过不断的试验找到了最优配置。其次，为了让气垫船能够在赛道上平稳前进并快速完成任务，我们采用了多级 PID 控制算法来调整船体姿态，同时通过卡尔曼滤波算法对传感器实时采集系统的位置信息进行数据处理，确保路径规划和导航的准确性。

该系统的设计和调试不仅让我提升了动手实践和编程调试的能力，也增强了我对多学科融合技术的理解与应用。竞赛和课程中的每一次挑战和收获，都是宝贵的经验。最重要的是，这次经历让我更加坚定了继续探索智能控制技术和创新设计的决心。
\subsection{改进展望}
在硬件层面，可以对无刷电机过零检测电路中的 C25、C26、C27 三颗滤波电容作移除处理。这三颗电容是用来滤除采集三点的高频分量的，但可能导致过零信号滞后，进而引发电机换相问题。

在软件层面，无刷电机启动直接采用了强拉换相方式。未来可以引入更为先进的正弦启动模式。通过使电机柔和地转动 360 度电角度后再切换至开环转动检测，可以极大地优化了启动过程，不仅能够降低启动冲击，还能显著提高了反电动势的检测精度，为电机的平稳运行奠定良好开端。此外，当前的多传感器数据融合算法主要基于卡尔曼滤波。未来可以探索更为复杂的融合方法，如扩展卡尔曼滤波（EKF）或无迹卡尔曼滤波（UKF），以提升姿态解算的精度和鲁棒性，尤其是在动态环境下的表现。